\chapter{Eosinophilic Granulomatosis with Polyangiitis}
\index{Eosinophilic Granulomatosis with Polyangiitis}
\label{sec:Eosinophilic Granulomatosis with Polyangiitis (EGPA)}

\section{Overview}
Eosinophilic granulomatosis with polyangiitis is a rare ANCA-associated vasculitis characterized by asthma, eosinophilia, and necrotizing vasculitis of small and medium-sized vessels with extravascular eosinophilic granulomas, with an incidence of 1--3 per million, equal sex distribution, and typically presenting in the fourth to fifth decade. The condition involves a Th2-driven immune response with prominent eosinophil activation and degranulation, ANCA (typically MPO-ANCA, positive in 30--50\%, more commonly in patients with kidney and nervous system involvement), and eosinophilic infiltration of tissues.
\section{Causes}
\section{Risk Factors}

\begin{itemize}[noitemsep]
  \item Genetic predisposition and family history
  \item Advancing age
  \item Lifestyle factors (diet, exercise, smoking, alcohol)
  \item Environmental exposures
  \item Underlying medical conditions
  \item Medication use
\end{itemize}
\section{Signs and Symptoms}

\subsection{Common symptoms}
\begin{itemize}[noitemsep]
  \item Pain or discomfort in the affected area
  \end{itemize}

\subsection{Emergency warning signs}
\begin{itemize}[noitemsep]
  \item Severe or worsening symptoms of eosinophilic granulomatosis with polyangiitis require immediate medical evaluation
\end{itemize}
\section{Progression and Stages}
Clinical features evolve through three phases: prodromal (asthma, allergic rhinitis, often years before systemic disease), eosinophilic (peripheral eosinophilia $>$1500/$\mu$L, eosinophilic pneumonia, gastroenteritis), and vasculitic (systemic vasculitis with mononeuritis multiplex---the most common vasculitic manifestation, purple spots on the skin, glomerulonephritis, cardiomyopathy, and skin nodules), with cardiac involvement (eosinophilic myocarditis, coronary arteritis, heart failure) being the leading cause of death. The condition may progress through different stages of severity over time. Early stages often involve mild or subtle symptoms that may be overlooked. Moderate stages involve more noticeable symptoms affecting daily function. Advanced stages may involve significant impairment and complications.
\section{Complications}
\begin{itemize}[noitemsep]
  \item Disease progression leading to worsening symptoms
  \item Secondary infections or injuries
  \item Side effects of treatment
  \item Reduced quality of life
  \item Emotional and psychological impact
\end{itemize}
\section{Diagnosis}
Doctors diagnose this based on the 2022 ACR/EULAR criteria. Comprehensive medical history and physical examination Laboratory tests to assess organ function and rule out other conditions Imaging studies to visualize affected structures Specialized testing depending on the affected system Consultation with relevant specialists Monitoring of disease progression over time

\begin{itemize}[noitemsep]
  \item Comprehensive medical history and physical examination
  \item Laboratory tests to assess organ function
  \item Imaging studies to visualize affected structures
  \item Specialized testing depending on the affected system
  \item Consultation with relevant specialists
\end{itemize}
\section{Prevention}
\begin{itemize}[noitemsep]
  \item Maintain a healthy lifestyle with balanced diet and exercise
  \item Avoid known risk factors
  \item Attend regular medical check-ups for early detection
  \item Follow recommended screening guidelines
\end{itemize}
\section{Treatment Options}
\begin{itemize}[noitemsep]
  \item Treatment typically involves stratified: patients without poor prognostic factors (FFS five-factor score=0) receive corticosteroids alone
  \item Patients with poor prognostic factors (FFS $\geq$1) receive corticosteroids plus cyclophosphamide or rituximab.
  \item Medications to manage symptoms and address underlying mechanisms
  \item Lifestyle modifications and self-care measures
  \item Physical therapy and rehabilitation
  \item Surgical intervention when indicated
  \item Regular monitoring and follow-up care
\end{itemize}
\section{Living with It}
\begin{itemize}[noitemsep]
  \item Follow your treatment plan as prescribed
  \item Maintain regular follow-up with your healthcare provider
  \item Make necessary lifestyle adjustments
  \item Seek support from family, friends, or support groups
  \item Stay informed about your condition
\end{itemize}
\section{Outlook}
\begin{itemize}[noitemsep]
  \item Your outlook depends on cardiac involvement and kidney function
  \item 5-year survival exceeds 90\%.
\end{itemize}
